% Ad Atticum 10.14 (ad-atticum-10-14)
% Source: https://raw.githubusercontent.com/PerseusDL/canonical-latinLit/master/data/phi0474/phi057/phi0474.phi057.perseus-lat1.xml
% Fetched by scripts/fetch_latin.py

% Dateline (Perseus): Scr. in Cumano viq i Id. Mai. a. 705 (49).

\ciceroLetterOpener{CICERO ATTICO salutem}

\ciceroSection{1} o vitam miseram maiusque malum tam diu timere quam est illud ipsum quod timetur! Servius, ut antea scripsi, cum venisset Nonis Maiis , postridie ad me mane venit. ne diutius te teneam, nullius consili exitum invenimus. numquam vidi hominem perturbatiorem metu; neque hercule quicquam timebat quod non esset timendum; illum sibi iratum, hunc non amicum; horribilem utriusque victoriam cum propter alterius crudelitatem, alterius audaciam, tum propter utriusque difficultatem pecuniariam; quae erui nusquam nisi ex privatorum bonis posset. atque haec ita multis cum lacrimis loquebatur ut ego mirarer eas tam diuturna miseria non exaruisse. mihi quidem etiam lippitudo haec, propter quam non ipse ad te scribo, sine ulla lacrima est sed saepius odiosa est propter vigilias.

\ciceroSection{2} quam ob rem quicquid habes ad consolandum conlige et illa scribe non ex doctrina neque ex libris (nam id quidem domi est, sed nescio quo modo imbecillior est medicina quam morbus), haec potius conquire de Hispaniis, de Massilia; quae quidem satis bella Servius adfert; qui etiam de duabus legionibus luculentos auctores esse dicebat. haec igitur si habebis et talia. et quidem paucis diebus aliquid audiri necesse est.

\ciceroSection{3} sed redeo ad Servium. distulimus omnino sermonem in posterum, sed tardus ad exeundum multo se in suo lectulo malle, quicquid foret. odiosus scrupulus de fili militia Brundisina. unum illud firmissime adseverabat, si damnati restituerentur, in exsilium se iturum. nos autem ad haec et id ipsum certo fore et quae iam fierent non esse leviora, multaque conligebamus. verum ea non animum eius augebant sed timorem, ut iam celandus magis de nostro consilio quam ad id adhibendus videretur. qua re in hoc non multum est. nos a te admoniti de Caelio cogitabimus.
